\section{Stage 248: Dynamic event-chain compiler}
\label{sec:app-part08:dynamic-event-chain-compiler}
\label{stage:248}
\stagefield{Verification}{SymPy audit: \StageFile{scripts/moving_throat_pde_stage248_dynamic_event_chain_compiler_from_relaxed_stationary_barrier_front_end_turning_point_threshold_speed_and_wkb_sympy_audit.py}.  Mathematica audit: \StageFile{mathematica/moving_throat_pde_stage248_dynamic_event_chain_compiler_from_relaxed_stationary_barrier_front_end_turning_point_threshold_speed_and_wkb_mathematica_audit.wl}.}


Stage 248 promotes \(V_{\rm eff}^{(247)}\) into a reduced scattering/event-chain compiler.  Write
\begin{equation}
\label{eq:app-part08-stage248-vdef}
V(r):=V_{\rm eff}^{(247)}(r).
\end{equation}
The reduced radial equation is
\begin{equation}
\label{eq:app-part08-stage248-radial-eom}
m_s\ddot r=-V'(r),
\end{equation}
with conserved energy \cref{eq:app-part08-intro-event-energy}.  If the branch has a single barrier peak
\begin{equation}
\label{eq:app-part08-stage248-peak}
V'(r_{\rm peak})=0,
\qquad
V''(r_{\rm peak})<0,
\qquad
V_{\rm peak}=V(r_{\rm peak}),
\end{equation}
then the finite-radius threshold law is \cref{eq:app-part08-main-vcrit-new}.  The pure-Coulomb contact threshold is \cref{eq:app-part08-main-vcrit-coul}.  The classical comparison window is therefore
\begin{equation}
\label{eq:app-part08-stage248-classical-window}
v_{\rm crit,new}<v_0<v_{\rm contact,Coul}.
\end{equation}

For subbarrier energy \(V(r_0)<E<V_{\rm peak}\), turning points solve
\begin{equation}
\label{eq:app-part08-stage248-turning-points}
V(r_-(E))=E,
\qquad
V(r_+(E))=E,
\qquad
r_-(E)<r_+(E),
\end{equation}
and transport as
\begin{equation}
\label{eq:app-part08-stage248-turning-transport}
\frac{dr_\pm}{dE}=\frac1{V'(r_\pm(E))}.
\end{equation}
The WKB action and transmission ratio are \cref{eq:app-part08-main-wkb-action,eq:app-part08-main-transmission-ratio}.  The pure-Coulomb outer turning point is exact,
\begin{equation}
\label{eq:app-part08-stage248-coul-turn}
r_{\rm turn,Coul}(E)=\frac1E,
\end{equation}
and the Coulomb action closes as
\begin{equation}
\label{eq:app-part08-stage248-coul-action}
I_{\rm Coul}(E)
=
\frac{\sqrt{2m_s}}{\hbar_{\rm eff}}
\left[
\frac{\pi}{2\sqrt E}
-
\sqrt{r_{\rm contact}(1-Er_{\rm contact})}
-
\frac{\arcsin\sqrt{Er_{\rm contact}}}{\sqrt E}
\right]
\end{equation}
for \(Er_{\rm contact}<1\).

Near the top, with
\begin{equation}
\label{eq:app-part08-stage248-near-top}
V(r)=V_{\rm peak}-\frac{K_{\rm peak}}2(r-r_{\rm peak})^2+O((r-r_{\rm peak})^3),
\qquad
K_{\rm peak}=-V''(r_{\rm peak}),
\end{equation}
the leading top action is
\begin{equation}
\label{eq:app-part08-stage248-top-action}
I_{\rm top}(E)=
\frac{\pi(V_{\rm peak}-E)}{\hbar_{\rm eff}}
\sqrt{\frac{m_s}{K_{\rm peak}}}
+O((V_{\rm peak}-E)^{3/2}).
\end{equation}
The dynamic diagnostics carried forward are
\begin{equation}
\label{eq:app-part08-stage248-diagnostics}
\Xi_{\rm turn}(E)=\Xi_1(r_+(E)),
\qquad
\lambda_{\rm th}(E)=\left|\frac{E}{V'(r_+(E))}\right|.
\end{equation}
These tags are later used in the helicity and Goldilocks compilers; they do not alter the scalar event chain.
